\note{sa page 7}

On a aussi l'\uncertain{homom.}\ $\Lambda$ de degré $-2$, défini par
(6.5) \ldots

\textbf{Théorème.} Les conditions suivantes sont équivalentes \add{(i) à (iii)} :

\begin{enumerate}
  \item Les conditions $(6.1)_{i}$ et $(6.2)_{i}$ pour tout $i$, et les $\pi_{i} \in \mathcal{E}$.
  \item \struck{\ill{}} donc la condition $(6.3)_{i}$ pour tout $i$.
  \item \struck{\ill{}} \add{$(6.1)_{i}$ pour tout $i$ (définissant $\Lambda$) et}
        $\Lambda \in \mathcal{E}^{(-2)}$ \quad (6.6)
  \item Les $\pi_{i,\alpha}$ sont dans $\mathcal{E}$ \ill{}
  \item Les $\pi_{i} \in \mathcal{E}$, \add{et ils sont} \uncertain{la seule famille} d'éléments de
        $\mathcal{E}$ satisfaisant (1.1.) et (1.3) a) et b), ou \ill{} (1.4).
  \item L'élément $\Lambda$ de $\mathcal{E}^{(-2)}$ est \add{uniquement} déterminé
        \struck{\ill{}} (en termes de \add{l'anneau gradué} $\mathcal{E}^{*}$ et de $L \in \mathcal{E}^{(2)}$),
        par les conditions suivantes :
\end{enumerate}

\note{La liste est très raturée : les énoncés (ii) à (v) sont écrits, biffés,
réécrits et renumérotés en marge, et seule leur charpente se lit. Ce qui est
sûr est que la chaîne va de la condition de Lefschetz (6.1) à l'existence dans
$\mathcal{E}$ de l'opérateur $\Lambda$ de degré $-2$, puis des projecteurs
$\pi_{i,\alpha}$ de la décomposition primitive.}

\[ (6.7) \qquad \begin{cases} L \Lambda = 1 - \sum_{0 \leqslant i \leqslant n} \pi_{i,0} \\[2pt] \Lambda L = 1 - \sum_{n \leqslant i \leqslant 2n} \pi_{i,0} \end{cases} \]

\note{Le (vi) et les relations (6.7) sont traversés en entier de grands traits
obliques, d'un trait plus fin et plus pâle que l'écriture : barrés, semble-t-il,
comme le reste de la liste. La page 13 reprend $\Lambda$ sous le numéro (6.15).}

\emph{Dém.} L'équivalence de (i) et (iii) est claire, et le fait que
(i) implique (iii) (en prenant $w_{i} = \Lambda^{n-i}$). % exponent: a figure before n is struck; w_i of degree -(2n-2i) forces Lambda^{n-i}
\uncertain{Prouvons} que (i), (ii) impliquent \ill{}. On \struck{\ill{}}
\add{\uncertain{a déjà indiqué}} que ça implique (v).

